\section{Related Work}

\subsection{Privacy-Preserving and Ultra-Low-Resolution Thermal Sensing}

Wearable indoor navigation systems for visually impaired users have traditionally relied on RGB cameras or depth sensors, which raises privacy concerns when deployed continuously indoors. Gao~\textit{et al.}'s DareDevil~\cite{gao2025daredevil} addresses this by auto-labeling thermal frames from a paired RGB camera through a pretrained YOLOv8 detector, avoiding manual annotation. However, DareDevil collapses all object classes into a single binary obstacle label and operates on colorized thermal frames from a single standard-resolution thermal camera at inference time; it does not address the extreme spatial compression of a $32\times24$-pixel array, multi-camera fusion, or robustness to sensor placement.

A separate line of work uses ultra-low-resolution thermal sensors (e.g., $24\times32$ arrays, on the order of \$30 per unit) as a low-power, lighting-invariant alternative to RGB and depth sensing. Such sensors are too coarse to resolve identifiable facial features, but their reduced spatial detail also degrades the features downstream perception models rely on; prior work compensates for this by fusing thermal imagery with gyroscope readings for motion estimation~\cite{thermal_gyro_odom}. Naser~\textit{et al.}~\cite{naser2022calibration} separately showed that most low-resolution thermal pipelines still assume a fixed, calibrated sensor geometry. This work targets discrete navigability classification rather than continuous odometry, and removes the dependence on geometric calibration that these systems retain.

\subsection{Automatic Label Generation for Traversability}

Manually annotating walkable regions is expensive and does not scale to real-world indoor environments. Kim~\textit{et al.}~\cite{kim2024traversability} addressed this for egocentric video with an automated annotation strategy that derives traversability labels without manual pixel-level annotation. We apply this same philosophy to a cross-modal setting: RGB-derived object detections supervise a thermal-only model only during data collection, and the RGB sensor is entirely absent at deployment.

\subsection{Large-Scale and Web-Sourced Navigation Learning}

CityWalker~\cite{liu2025citywalker} trains an embodied urban navigation policy on over 2{,}000 hours of in-the-wild web video, using visual odometry to generate pseudo-labels instead of teleoperated demonstrations. CityWalker targets continuous waypoint-following from RGB video; this work targets discrete per-direction navigability from privacy-preserving thermal sensing, with an order of magnitude less data collected from a purpose-built wearable rig.

\subsection{Multi-View Sensor Fusion and Placement Robustness}

Multi-camera perception systems typically assume a fixed mapping between physical sensor position and input channel, an assumption that breaks when sensors shift or are reattached in a different order on a wearable rig. Pixel-level, decision-level, and feature-level fusion are well studied in the sensor fusion literature, but little work examines how these strategies behave when the camera-to-slot assignment is perturbed, particularly at the small camera counts and extreme resolutions considered here. This also relates to cross-resolution knowledge transfer, where a high-resolution signal supervises a model on coarser inputs~\cite{qi2021msad}; our gap is larger still, since the supervisory RGB detector and the deployed $32\times24$ thermal arrays share no feature space and no geometric alignment at inference time. We contribute a direct comparison of three fusion strategies for an ultra-low-resolution, three-camera array, together with an explicit treatment of camera-slot permutation robustness.
